\section{Summarization Techniques}\label{sec:summarization}
In this section, we review the summarization techniques utilized in our application. As it's mentioned in section \ref{sec:technology}, we already had an application capable of creating summaries. The main contribution made by our group is extending the summary capabilities with topic-aware summarization and delta summarization. \\
To have a basic understanding of how summarization is made in our application, we should start explaining the existing model and how we utilized it. Although the summarization model can be configured via a configuration file, we used "philschmid/bart-large-cnn-samsum" from Hugging Face (\cite{schmid2021bart}). This model is pre-trained with the samsum dataset (\cite{samsum}). Since it contains several conversations and their respective summaries, the pre-trained model matched our objectives. We knew that BART has some limitations such as the token limit of 1024 but decided to use it as the main focus of this project is integrating the topic-aware summarization and delta summarization modules. To overcome the maximum token limit problem, we leverage chunking to split the text into smaller pieces. By having different text pieces with smaller sizes, our BART model is capable of producing a summary from those. In other words, the overall summary consists of the summaries generated from the chunked parts.
\subsection{Topic-Aware Summarization}
Topic modeling was one of the already implemented features in our app. Therefore, we could extract the topics from a given set of documents. Our first approach is using the conversations streamed from Reddit and Telegram and feed them as documents to our BERTopic \footnote{\href{https://maartengr.github.io/BERTopic/index.html}{BERTopic, https://maartengr.github.io/BERTopic/index.html}} model. Since we did not have the API implementations in the first place, the Samsum dataset was used for the conversation data. Topic-aware summarization flow can be described as the following:
\begin{itemize}
    \item Extract topics and respective topic embeddings from the streamed conversations utilizing BERTopic.
    \item Extract sentences from the dialogue to be summarized with respect to a similarity score. Here, we use a sentence transformer, "paraphrase-multilingual-MiniLM-L12-v2" (\cite{reimers-2019-sentence-bert}), from Hugging Face to create the embeddings of the sentences. By comparing the sentence embeddings of the dialogue with the topic embeddings, a similarity matrix is generated. Each sentence is mapped to the most similar topic, and, eventually, for each topic (potentially) we have some extracted sentences.
    \item Finally, we have extracted sentences for each topic and as a final layer, they feed into our BART summarizer model for abstractive summarization.
\end{itemize}
One major challenge we face is the lack of conversations fetched from the APIs. Since our app is available in English, we filtered out the conversations in different languages. This created a bottleneck in the usage of BERTopic because it requires a minimum number of documents to generate a meaningful output. As a fallback mechanism, we introduced static pre-defined topics. If BERTopic fails to create the topics, we use static topics and their embeddings to create the topic-aware summarizations.
\subsection{Delta Summarization}
Delta summarization may sound misleading in the context of our application as we do not create a summary but a comparison. The main idea is to compare two different topic summaries and visualize them in a certain format. The currently supported formats are cosine similarity matrices and scatter plots (2D).
\subsubsection{Cosine Similarity Matrix}
It's a matrix where the items are cosine similarity values of the two respective topic summaries. We first find the cosine similarity value between two summaries by using their sentence embeddings. The sentence encoder, "all-MiniLM-L6-v2", is utilized to create those embeddings. It's a pre-trained model from Hugging Face and is from the sentence-transformers family. After having the cosine similarity values, we create the matrix as an output (\ref{fig:delta-summary-cosine}).
\subsubsection{Scatter Plot}
To create this format, we follow a common approach. That is, we create the sentence embeddings as it's been done in the cosine similarity matrix. Having this, Principal Component Analysis (PCA) is applied for dimensionality reduction to the high-dimensional embeddings. In other words, sentence embeddings created by the sentence transformer have 384 dimensions and to be able to show it in a 2D plot, we must apply a dimensionality reduction technique. PCA is selected because it's easily applicable from open-source libraries and is also a well-known method. After all, we see a 2D scatter plot showing how similar the two summaries are (\ref{fig:delta-summary-scatter}). This is helpful since we can observe the similarity of different topic summaries and also compare them with the original dialogue and the default summary.